problem:  
Let \((M^6,J,g)\) be a compact, simply connected **strict nearly Kähler** manifold with SU(3)-structure \((\omega,\psi_+,\psi_-)\), satisfying  
\[
d\omega = 3\,\psi_+ , \quad d\psi_- = -2\,\omega^2.
\]
Let \(\nabla^c\) be the canonical Hermitian connection with totally skew-symmetric torsion \(T = \tfrac{1}{2}d\omega\), and let \(|W_1|\) be the constant norm of its intrinsic torsion (normalized so that \(\mathrm{Ric}(g) = 5 |W_1|^2 g\)).  

From standard Weitzenböck-type formulas on Hermitian manifolds with parallel torsion, primitive \((1,1)\)-forms \(\alpha \in \Lambda^{1,1}_0\) satisfy schematically
\[
\Delta_g \alpha = (\nabla^c)^* \nabla^c \alpha + Q(T)(\alpha),
\]
where \(Q(T)\) is an algebraic curvature–torsion operator determined by contractions of \(T\) and \(J\).  
For homogeneous strict nearly Kähler examples—\(S^6\), \(S^3 \times S^3\), \(\mathbb{C}P^3\), and the flag manifold \(F^3\)—representation-theoretic computations indicate that the first nonzero eigenvalue \(\lambda_1^{(1,1)}\) of the Hodge Laplacian on primitive \((1,1)\)-forms is proportional to \(|W_1|^2\).

**Problem:**  
Formulate and test the following **Spectral Rigidity Conjecture for Nearly Kähler 6‑manifolds**:

1. Prove or disprove the existence of a universal constant \(C_0>0\) (depending only on normalization of \(|W_1|\)) such that for every compact strict nearly Kähler 6‑manifold,
   \[
   \lambda_1^{(1,1)} \ge C_0\,|W_1|^2.
   \]
   The conjectural value \(C_0\) should coincide with that computed on known homogeneous examples.
2. Determine whether **equality** in this inequality characterizes the *homogeneous* nearly Kähler 6‑manifolds.
3. More generally, investigate whether the quantity \(\lambda_1^{(1,1)}/|W_1|^2\) is a *rigid invariant* that distinguishes homogeneous from inhomogeneous nearly Kähler metrics.

Approaches might include deriving an explicit Weitzenböck formula for primitive \((1,1)\)-forms under \(\nabla^c\), analyzing curvature–torsion terms that govern the spectrum of \(\Delta_g\), or constructing numerical/analytic comparisons with the known homogeneous models.

---

Why is it a "good" problem:  
This version refines the question into a precisely stated **spectral–analytic rigidity conjecture** rooted in the geometry of nearly Kähler manifolds. The inequality and rigidity test are well-defined, compatible with known normalization conventions, and grounded in the intrinsic torsion structure.  

- It is **profound**: it seeks a universal curvature–torsion constraint linking the intrinsic torsion constant to the first Laplacian eigenvalue—a conjectural “Bochner inequality with torsion.”  
- It is **bridging**: the problem connects **special holonomy theory (SU(3) with torsion)** and **spectral geometry**, possibly enabling classification or detection of homogeneity through analytical invariants.  
- It is **simple yet deep**: the question “does the first eigenvalue of the Laplacian on primitive (1,1)-forms control homogeneity of strict nearly Kähler manifolds?” is clear, self-contained, and potentially decisive in elucidating analytic rigidity phenomena in 6‑dimensional special geometry.  

If resolved, it would reveal whether homogeneous nearly Kähler manifolds are characterized by optimal spectral bounds among all strict nearly Kähler structures—opening an analytic path toward geometric classification beyond existing algebraic methods.